\section{Joint response inference, persistence and learned adjustment}
\label{s4:response-results}
The response experiments instantiate Section~\ref{s4:response-construction}
with two banking outcomes, 13 horizons from zero through twelve, and a
common set of origins. Each row of Table~\ref{s4:response-coverage} contains
500 independent replications. Whole-curve coverage requires all 26 true
coordinates to lie in their reported sets. Applying 95\% pointwise intervals
to every plotted horizon is therefore deliberately an adverse comparison.
Zero numerical failures occurred; failed replications, had they occurred,
would remain in the denominator.

\begin{table}[htbp]\centering\small
\begin{tabular}{lrrrrrr}
\toprule
Case & Iid & HAC & Joint & Max & IV set & VAR\\
\midrule
Strong, white & 0.434 & 0.414 & 0.958 & 0.930 & 0.952 & 0.950\\
Strong, persistent & 0.132 & 0.452 & 0.934 & 0.886 & 0.930 & 0.938\\
Persistent, long & 0.116 & 0.468 & 0.926 & 0.876 & 0.928 & 0.978\\
Weak, persistent & 0.670 & 0.850 & 0.962 & 0.944 & 0.924 & 0.986\\
Near unit & 0.246 & 0.598 & 0.818 & 0.724 & 0.776 & 0.966\\
Invalid IV & 0.002 & 0.012 & 0.182 & 0.136 & 0.270 & 0.000\\
\bottomrule
\end{tabular}

\caption{Whole-curve coverage of nominal 95\% procedures. ``Joint'' uses
Bonferroni; ``max'' uses the estimated Gaussian maximum; ``IV set'' inverts
the joint residual tests. Standard Monte Carlo errors near .95 are about
.010. These are synthetic experiments, not coverage guarantees.}
\label{s4:response-coverage}
\end{table}

Strong instruments have innovation loading one; the weak case has loading
$.06$. The reference system has spectral radius $.81$; the near-unit case
uses $.985$. Persistent instrument noise has AR coefficient $.85$, rising
to $.95$ in the near-unit case. Most experiments use 800 origins; the long
persistent case uses 5,000. The invalid instrument additionally loads on
another structural innovation by $.6$. The initial stationary law and
all generating matrices are fixed in the reproducible specification.

Several method decisions follow. Joint covariance matters because a
pointwise interval is not a curve guarantee. Yet a consistent HAC estimator
can still underestimate finite-sample uncertainty: the Gaussian-maximum
band covers only .886 and .876 in the two persistent strong-instrument
experiments, and .724 near the unit boundary. Bonferroni is more conservative,
but its .818 near-unit coverage also fails the preset adequacy threshold.
These failures prohibit presenting the default bandwidth as automatically
calibrated for persistent bank series. The statistical result remains
asymptotic under a fixed stability margin, not uniform near a unit root.

The proxy VAR has much smaller mean squared error in the correctly
specified finite-order generating law. That advantage is a consequence of
its propagation restriction. It does not establish that an economist should
replace direct moments when finite-order dynamics are misspecified. Weak
normalization makes the direct and proxy mean squared errors about 752
and 13.5; 96.4\% of runs yield an unbounded residual-test set in at least
one coordinate. A narrow decision recommendation would conceal the
information deficiency. Instrument invalidity defeats both routes, including
zero proxy-VAR coverage in this example.

A separately specified follow-up, with a new random seed, compares bandwidth
$L=\lfloor n^{1/4}\rfloor$ and $4L$ at 5,000 origins. Their whole-curve
coverages are .948 and .968, with mean widths .312 and .338. This new
experiment does not replace the original .926 finding or establish that
$4L$ is universally preferable. The selected workflow reports sensitivity
to the lag window, compares the structural challenger when justified and
withholds a calibrated finite-sample claim in regimes that fail the stated
qualification. Picking the more favorable seed or interval after observing
the answer would be a different, unqualified procedure.

The integrated experiment has one common random-walk trend and two stationary
spreads. Mapping stationary-change IV estimates to cumulative level responses
gives whole-curve coverage .966 and mean squared error .0479. Treating the
levels as stationary gives .586 coverage and error 55.66. Cointegration is
substantive here: it defines which components are stationary and which
intervention response is being summed. The comparison does not certify
unknown rank selection or arbitrary local-to-unit triangular arrays.

With 4,000 dates and 80 bounded controls, separated-block sparse IV gives
.942 coverage of the scalar effect $1.2$, compared with .950 for oracle
nuisances. A three-factor reduction and an unadjusted estimator both have
zero coverage, with mean squared errors about .058 versus .00102 for
sparse adjustment. The causal confounder is sparse in this generator and
need not align with its leading covariance components. Factor compression
requires an information-preservation argument, not only a prediction score.
The fixed penalty multiplier in this experiment establishes feasibility;
the required population nuisance rates remain assumptions to justify in
each application.

For a fixed 96-bank panel, the additional mean-score experiment separates
common and bank-specific Gaussian AR components. At 800 dates, row-iid
intervals cover in .116 of runs; date-aggregate HAC covers in .856; an
oracle using the exact finite-sample Gaussian variance covers in .938.
At 2,000 dates the corresponding values are .104, .890 and .946.
Correctly recognizing the date as the information unit removes a gross
variance error, but the short Bartlett lag window still underestimates
persistent uncertainty. More customer rows do not repair either problem.

\section{Complete computational workloads and their statistical meaning}
\label{s4:workflow-results}
Table~\ref{s4:workflow-cost} measures fresh processes, including imports,
synthetic generation, local Parquet write/read, observation reconstruction,
fitting, uncertainty, the specified comparison or resampling, and report
serialization. Stage timings, physical and logical I/O, peak resident
memory, seeds, numerical differences and optimizer stops are retained.
Each row is the median of three processes on an AMD EPYC 7773X host with
256 logical CPUs; BLAS, DuckDB and Polars are each restricted to one thread.
The machine and library versions are recorded with the results. The
measurements are local warm-filesystem workloads, not network/object-store
or distributed-system benchmarks.

\begin{table}[htbp]\centering\small
\begin{tabular}{lrrr}
\toprule
Workload & Raw rows & Seconds & Peak MiB\\
\midrule
Response, small & 77,952 & 2.00 & 267.7\\
Response, medium & 772,608 & 2.22 & 356.8\\
Response, large & 3,849,216 & 3.38 & 717.2\\
Learning, small & 4,000 & 1.56 & 256.0\\
Learning, large & 24,000 & 2.07 & 432.0\\
Latent, small & 3,200 & 7.49 & 254.0\\
Latent, large & 38,400 & 48.27 & 264.0\\
Temporal graph, small & 1,000 & 1.49 & 250.0\\
Temporal graph, large & 4,000 & 1.71 & 257.2\\
\bottomrule
\end{tabular}

\caption{Complete synthetic workload cost. Response rows include three raw
versions per bank/date; learning rows include two versions per date.
Latent rows include 19 whole-trajectory bootstrap refits with three starts.
Memory includes the interpreter, libraries and retained synthetic arrays.}
\label{s4:workflow-cost}
\end{table}

For the response workload, bank deviations are centered across banks so
that the fixed weighted aggregate equals the common structural series.
Raw data contain an initial publication, an eligible correction and a
future revision. The SQL operation restricts availability before selecting
the highest eligible version, then aggregates to dates. It is checked
against the separately retained generating series, and the fitted curve
against its direct numerical reference. Maximum discrepancies are below
$10^{-9}$ at every scale. The largest case reconstructs 3,849,216 versioned
records for 256 banks and 5,012 dates, then fits direct and proxy responses,
joint covariance, residual-test sets and a bandwidth comparison. It is
one common-shock analysis; the record count does not create millions of
independent macroeconomic shocks.

The learned-control workload reconstructs a wide versioned table and fits
both separated-block sparse adjustment and the factor challenger, retaining
their errors and covariance. The large case has 12,000 dates and 200 controls.
The temporal graph workload fits 12 variables at 4,000 dates and propagates
the joint coefficient box through 13 horizons. Its timing does not establish
uniform high-dimensional graph coverage; the inferential reference remains
the declared fixed-dimensional structural class. Exponential DAG enumeration
and the 25-DAG Gaussian reference are different workloads.

The latent workload reconstructs complete customer trajectories, fits the
unknown-emission model with multiple starts, computes observed-information
value uncertainty and refits every bootstrap sample. The larger case has
3,200 customers and 12 decisions. Its model fit and information calculation
cost only a small fraction of the 19-refit workflow. Nineteen bootstrap
draws are suitable for demonstrating cost scaling, not precise 95\% interval
endpoints. The separate 79-refit fixed-policy example and repeated-sampling
study address different inferential questions. All caps, boundary and
optimizer diagnostics remain in the receipts, including intermediate EM
caps followed by likelihood refinement.

These results change the engineering priority. For short response curves,
data reconstruction and process overhead can dominate the numerical moment
calculation. For controlled likelihoods, repeated fitting dominates.
Caching immutable views helps the first workload; efficient likelihood
scores, warm initial proposals checked against independent starts and
parallel independent bootstrap fits can help the second. Each proposed
acceleration must retain the same target and numerical tolerance. No
bank-level labor saving follows from these CPU timings: the expert
composition comparator can use the same engines and estimators.
